A Seção~\ref{sec:antispoiler} garante que o assistente nunca exponha o que o leitor ainda não alcançou, e os modelos geradores decidem \textit{o que} o assistente responde a cada solicitação; nenhum dos dois, porém, estabelece se uma resposta gerada está de fato \textit{correta}, nem comunica ao leitor o quanto o sistema confia nela. É essa lacuna que a camada de validação preenche. Atribuí-la ao sistema não é um detalhe de implementação: a ``comunicação cuidadosa das respostas, exibindo limitações e incerteza'' figura entre os requisitos \textit{Must} do projeto (Tabela~1), e esta seção descreve o mecanismo --- o \textit{validador} --- com que o realizamos.

O ponto de partida da camada de validação é o reconhecimento de que ``correto'' não é uma propriedade única. Uma resposta gerada pelo assistente costuma reunir afirmações de naturezas distintas: um fato de enredo ancorado no texto já lido, uma paráfrase que precisa preservar o sentido do trecho original, um sentido de palavra que deve corresponder ao registrado em dicionário, ou ainda um fato sobre o mundo real que vive fora da obra. Cada uma dessas naturezas se verifica de forma diferente e, sobretudo, contra uma fonte de evidência diferente. Submeter todas a um único \textit{prompt} do tipo ``isto está correto?'' herdaria justamente o pior comportamento desse tipo de instrução: o erro confiante, em que o modelo afirma com a mesma segurança um fato que recuperou do texto e um que apenas recordou de memória --- uma alucinação. É essa indistinção que a validação precisa desfazer, tratando cada tipo de afirmação com a metodologia de checagem mais apropriada a ele.

Essa correspondência entre tipo de afirmação e checagem se concretiza de modo diferente conforme a intenção que o leitor aciona (ver Seção~\ref{sec:visao-geral-sistema}). Para contextualizar e relembrar, o validador decompõe a resposta gerada em afirmações menores e auto-contidas --- resolvendo pronomes e separando fatos de atribuições causais --- e verifica cada uma isoladamente contra o contexto recuperado, em um \textit{prompt} próprio, de modo que receba seu próprio veredito. A funcionalidade de definição segue uma estrutura diferente: em vez de decompor a resposta, o validador trata a definição geral do trecho como uma unidade, checando-a contra um contexto recuperado reduzido, e, em \textit{prompts} separados, valida a definição de cada palavra difícil sinalizada (Seção~\ref{sec:dificuldade-palavras}); esta última etapa apoia-se em um dicionário \textit{offline} (WordNet, via \textit{nltk} \cite{miller1992wordnet}) como fonte de evidência lexical. A paráfrase, por fim, é entregue ao validador de forma integral, sem decomposição e sem recuperação de contexto: o que se verifica é se a frase reescrita preserva o sentido do trecho original --- uma relação de equivalência que se perderia caso a frase fosse fragmentada.
